\section*{Chapter \ref{ch:qm:transforms-interpolation-and-algorithmic-differentiation} --- Transforms, Interpolation and Algorithmic Differentiation}

\begin{solution}{exo:qm:transforms-interpolation-and-algorithmic-differentiation:1}
Directly $N^2 = 16\,777\,216$; radix 2 about $\frac N2\log_2N = 2\,048 \times 12 = 24\,576$, some 700 times fewer.
\end{solution}

\begin{solution}{exo:qm:transforms-interpolation-and-algorithmic-differentiation:2}
$(1 + \epsilon)^2 = 1 + 2\epsilon$ and $e^{1+\epsilon} = e + e\epsilon$, so $f(1 + \epsilon) = (1 + 2\epsilon)(e + e\epsilon) = e + 3e\epsilon$: $f'(1) = 3e \approx 8.155$.
\end{solution}

\begin{solution}{exo:qm:transforms-interpolation-and-algorithmic-differentiation:3}
$5 \cdot 2^{-k} < 10^{-10}$ needs $k > \log_2(5 \times 10^{10}) = 35.5$: 36 steps, as in the tutorial.
\end{solution}

\begin{solution}{exo:qm:transforms-interpolation-and-algorithmic-differentiation:4}
\omterm{def:qm:transforms-interpolation-and-algorithmic-differentiation:tape}{Tape}: $v_1 = x_1$, $v_2 = x_2$, $v_3 = v_1v_2$ (partials $v_2$, $v_1$), $v_4 = \sin v_3$ (partial $\cos v_3$), $v_5 = v_4 + v_1$ (partials 1, 1). At $(1, 2)$: $v_3 = 2$. Reverse: $\bar v_5 = 1$; $\bar v_4 = 1$,
$\bar v_1 = 1$; $\bar v_3 = \cos2 = -0.416$; $\bar v_1 \mathrel{+}= \bar v_3v_2 = -0.832$, so $\bar v_1 = 0.168$; $\bar v_2 = \bar v_3v_1 = -0.416$. Check: $\partial y/\partial x_1 = x_2\cos(x_1x_2) + 1$, $\partial y/\partial x_2 = x_1\cos(x_1x_2)$.
\end{solution}

\begin{solution}{exo:qm:transforms-interpolation-and-algorithmic-differentiation:5}
$-\ln P(t) = r(t)t$, so $f(t) = \frac{d}{dt}(r(t)t) = r(t) + tr'(t)$. At the long end $t$ is large, so a spline slope $r'$ of $-0.001$ a year at $t = 30$ lowers the forward by 3 percentage points: small wiggles of the
zero curve become large wiggles of the forward, most of all at long maturities.
\end{solution}

\begin{solution}{exo:qm:transforms-interpolation-and-algorithmic-differentiation:6}
The payoff $\mathbf 1\{S_T > K\}$ is piecewise constant in the path, so its pathwise derivative is zero almost surely: differentiating each simulated payoff misses the jump. Remedies: smooth the payoff (a
narrow call spread), use the likelihood-ratio method for that term, integrate the last step analytically (\omterm{def:qm:probability-at-speed:condexp}{conditional expectation}), or combine them (Vibrato Monte Carlo).
\end{solution}

\begin{solution}{exo:qm:transforms-interpolation-and-algorithmic-differentiation:7}
\texttt{implied\_vol\_iterations()}: bisection 36 at every strike; secant 10, 7, 4, 7, 18 and failure; Newton 9, 6, 3, 6, 13 and failure; Brent 17, 17, 5, 17, 17 and 18. At 400 the price
($2.5 \times 10^{-7}$) and the vega at the starting guess ($8 \times 10^{-9}$) are tiny: Newton's first step overshoots to a volatility of 30.6, where the vega underflows, and the secant slope is equally
unreliable; the bracketing methods cannot leave $[0.001, 5]$.
\end{solution}

\begin{solution}{exo:qm:transforms-interpolation-and-algorithmic-differentiation:8}
A central difference has \omterm{def:qm:finite-difference-methods:consistent}{truncation error} of order $h^2f^{\prime\prime\prime}$ and rounding error of order $u|f|/h$; at the book's scale (value about 112, second and third derivatives large near the caplets' strikes) a
discrepancy of $10^{-5}$ on 500 is a relative $2 \times 10^{-8}$, what the bump itself is accurate to. The adjoint agrees with \omterm{def:qm:transforms-interpolation-and-algorithmic-differentiation:ad}{forward mode} to $10^{-13}$: the bump is the less accurate of the two.
\end{solution}

\begin{solution}{pb:qm:transforms-interpolation-and-algorithmic-differentiation:1}
\textbf{1.} Exponential decay: Black--Scholes exact to $7 \times 10^{-15}$ at 64 terms; Merton $10^{-7}$ at 64 and $2 \times 10^{-12}$ at 128.
\textbf{2.} $\P(S_T > 110) = 0.3503999$, against the Poisson mixture of lognormals to $3 \times 10^{-10}$.
\textbf{3.} Against 400\,000 simulated paths: 9.5243 against $9.528 \pm 0.022$.
\textbf{4.} The put payoff is bounded, so its cosine coefficients stay small and truncation of the upper tail does not amplify errors; parity gives the call.
\textbf{5.} When a whole strip of strikes is needed at once on a regular grid, as in Carr and Madan's method.
\textbf{6.} They fall to 2.63\% at thirty years: the natural end condition lets the spline slope down steeply, and $tr'(t)$ amplifies the slope at long maturities.
\textbf{7.} Monotone $rt$, hence nonnegative forwards whenever the data allow; here every forward lies between 3.32\% and 5.67\%.
\textbf{8.} Its forwards jump at every pillar, by up to 0.43 percentage points.
\textbf{9.} At the money Newton 3, secant 4, Brent 5, bisection 36; at 400 Newton and secant fail, Brent needs 18 and bisection 36.
\textbf{10.} 1.92 near the ends.
\textbf{11.} 2\,864 elementary operations and 3\,698 (operation, input) pairs.
\textbf{12.} Bumping 401 evaluations, \omterm{def:qm:transforms-interpolation-and-algorithmic-differentiation:ad}{forward mode} 1\,433, \omterm{def:qm:transforms-interpolation-and-algorithmic-differentiation:ad}{reverse mode} 3.58 by the operation count.
\textbf{13.} Reverse and forward agree to $2 \times 10^{-13}$; reverse and central bumps to $9 \times 10^{-6}$, the truncation and rounding error of the bumps.
\textbf{14.} With linear interpolation each cash-flow date depends on the two pillars around it; pillars with no date in their neighbourhood have zero sensitivity.
\textbf{15.} It holds at most 201 states instead of 10\,000 (60\,003 \omterm{def:qm:transforms-interpolation-and-algorithmic-differentiation:tape}{tape} nodes), for one extra forward pass, with the same gradient to $10^{-14}$.
\textbf{16.} Sensitivities of discontinuous payoffs and of inputs that enter through non-differentiable code (lookups, branches), and a nightly sample as a check.
\textbf{17.} Adjoint against bump on a sample of inputs, \omterm{def:qm:transforms-interpolation-and-algorithmic-differentiation:tape}{tape} size and memory, and that the \omterm{def:qm:transforms-interpolation-and-algorithmic-differentiation:tape}{tape} was cleared per valuation.
\textbf{18.} The real cost, including \omterm{def:qm:transforms-interpolation-and-algorithmic-differentiation:tape}{tape} memory traffic: operation counts ignore it; the C++ test measures under twenty evaluations.
\textbf{19.} Named result: \emph{one sweep instead of 401}: the adjoint gradient of the 400-input book costs 3.58 evaluations by operation count, against 401 for bumping, and agrees with the bumped gradient to
$9 \times 10^{-6}$ (with \omterm{def:qm:transforms-interpolation-and-algorithmic-differentiation:ad}{forward mode} to $2 \times 10^{-13}$).
\textbf{20.} Risk asks for many inputs' sensitivities of one number, and the \omterm{def:qm:transforms-interpolation-and-algorithmic-differentiation:ad}{reverse mode}'s cost does not depend on how many inputs there are.
\end{solution}

\begin{solution}{iq:qm:transforms-interpolation-and-algorithmic-differentiation:1}
\omterm{def:qm:transforms-interpolation-and-algorithmic-differentiation:ad}{Forward mode} carries derivatives along one input direction with the values: cost proportional to the number of inputs, cheap for few inputs and many outputs. \omterm{def:qm:transforms-interpolation-and-algorithmic-differentiation:ad}{Reverse mode} records the computation and
propagates adjoints of one output back to all inputs: cost a small multiple of one evaluation, independent of the number of inputs, cheap for one output and many inputs, at the price of memory.
\iqlookfor{Cost scaling with inputs versus outputs, and the memory of the \omterm{def:qm:transforms-interpolation-and-algorithmic-differentiation:tape}{tape}.}
\end{solution}

\begin{solution}{iq:qm:transforms-interpolation-and-algorithmic-differentiation:2}
An active type whose operators push (inputs, partials) onto a \omterm{def:qm:transforms-interpolation-and-algorithmic-differentiation:tape}{tape}, and a reverse loop over the \omterm{def:qm:transforms-interpolation-and-algorithmic-differentiation:tape}{tape} accumulating adjoints. Limits: memory of the \omterm{def:qm:transforms-interpolation-and-algorithmic-differentiation:tape}{tape} (\omterm{def:qm:transforms-interpolation-and-algorithmic-differentiation:tape}{checkpointing}), thread safety (a \omterm{def:qm:transforms-interpolation-and-algorithmic-differentiation:tape}{tape} per
thread), third-party code that is not templated, branches and kinks, and the overhead of taping inner loops, which expression templates or hand-written adjoints of kernels reduce.
\iqlookfor{\omterm{def:qm:transforms-interpolation-and-algorithmic-differentiation:tape}{Tape} design, memory and at least two production limits.}
\end{solution}

\begin{solution}{iq:qm:transforms-interpolation-and-algorithmic-differentiation:3}
Bracket the volatility (the price is increasing in it between known bounds), use a safeguarded method such as Brent's, or Newton with a fallback to bisection, starting from a good rational
approximation; work with normalised prices, and treat prices below intrinsic or above the upper bound as errors.
\iqlookfor{Bracketing with a fast interior step, and bounds checks.}
\end{solution}

\begin{solution}{iq:qm:transforms-interpolation-and-algorithmic-differentiation:4}
Interpolate something whose shape carries the economics, such as $\ln P(t)$ or $rt$, with a local monotone scheme, so that forwards stay positive, continuous where desired, and local: a quote change
moves the curve near its pillar only. A natural spline of zero rates is global (one quote moves all forwards) and can produce oscillating or negative forwards.
\iqlookfor{Locality and the forward curve.}
\end{solution}

\begin{solution}{iq:qm:transforms-interpolation-and-algorithmic-differentiation:5}
By Fourier methods: the COS expansion or Carr--Madan's damped transform with a \omterm{def:qm:transforms-interpolation-and-algorithmic-differentiation:dft}{fast Fourier transform}, or a Gil-Pelaez integral for probabilities; check the truncation range from the \omterm{def:qm:jump-processes:cumulant}{cumulants}, and
validate against simulation.
\iqlookfor{COS or FFT, with a truncation choice and a check.}
\end{solution}

\begin{solution}{iq:qm:transforms-interpolation-and-algorithmic-differentiation:6}
The barrier indicator is piecewise constant along each path, so the pathwise derivative misses the discontinuity. Use a smoothed barrier (a Brownian-bridge survival probability per step, which is
differentiable), a likelihood-ratio term, or \omterm{def:qm:probability-at-speed:condexp}{conditional expectations} at the barrier; then check against a bumped price with common random numbers.
\iqlookfor{Discontinuity and a differentiable reformulation.}
\end{solution}
